\chapter{Linear Algebra, Revisited}

\begin{center}\itshape
Why is $F_{n+1}F_{n-1}-F_n^2$ always $\pm1$? Why is the determinant of $\left(\frac1{i+j-1}\right)_{i,j}$ positive, though tiny?\\
Matrices answer questions that do not mention matrices.
\end{center}

\section{Determinants}

\begin{reminder}[determinants]
\begin{itemize}
\item $\det$ is the unique multilinear, alternating function of the columns with $\det I_n=1$.
\item $\det(AB)=\det A\det B$, $\det A^{\mathsf T}=\det A$, and $A$ is invertible if and only if $\det A\ne0$.
\item Adding to a row (or column) a multiple of another does not change the determinant; the determinant of a triangular matrix is the product of its diagonal.
\item \textbf{Rank--nullity}: for $u:E\to F$ linear with $\dim E<\infty$, $\dim\ker u+\operatorname{rk}u=\dim E$.
\end{itemize}
\end{reminder}

\begin{proposition}[Vandermonde]\label{prop:vdm}
For $x_1,\dots,x_n\in\mathbb K$,
\[
V(x_1,\dots,x_n)=\det\left(x_i^{\,j-1}\right)_{1\le i,j\le n}=\prod_{1\le i<j\le n}(x_j-x_i).
\]
\end{proposition}
\begin{proof}
Induction on $n$; $n=1$ is clear. If two of the $x_i$ are equal both sides vanish. Otherwise, replacing $x_n$ by a variable $X$ and expanding along the last row,
$P(X)=V(x_1,\dots,x_{n-1},X)$ is a polynomial of degree at most $n-1$ whose coefficient of $X^{n-1}$ is $V(x_1,\dots,x_{n-1})$. It vanishes at $x_1,\dots,x_{n-1}$ (two equal rows),
so $P=V(x_1,\dots,x_{n-1})\prod_{i<n}(X-x_i)$. Evaluate at $X=x_n$ and apply the induction hypothesis.
\end{proof}

(The circulant determinant of Exercise~\ref{ex:circulant} is another closed form, obtained by diagonalising.)

\section{Reduction}

\begin{reminder}[eigenvalues and diagonalisation]
Let $A\in M_n(\mathbb K)$ and $\chi_A(X)=\det(XI_n-A)$, of degree $n$, with $\chi_A=X^n-\operatorname{tr}(A)X^{n-1}+\dots+(-1)^n\det A$.
\begin{itemize}
\item The eigenvalues of $A$ are the roots of $\chi_A$ in $\mathbb K$. Over $\C$, $\operatorname{tr}A$ and $\det A$ are the sum and the product of the eigenvalues, counted with multiplicity.
\item $A$ is diagonalisable if and only if it is cancelled by a split polynomial with simple roots (equivalently, its minimal polynomial is split with simple roots).
\item \textbf{Cayley--Hamilton}: $\chi_A(A)=0$.
\item \textbf{Spectral theorem}: a real symmetric matrix $S$ has real eigenvalues and an orthonormal basis of eigenvectors: $S=PDP^{\mathsf T}$ with $P$ orthogonal and $D$ diagonal.
\end{itemize}
\end{reminder}

\begin{example}[linear recurrences]
If $u_{n+2}=au_{n+1}+bu_n$, then $\binom{u_{n+1}}{u_n}=M^n\binom{u_1}{u_0}$ with $M=\begin{pmatrix}a&b\\1&0\end{pmatrix}$, whose characteristic polynomial is $X^2-aX-b$.
When it has two distinct roots $r_1,r_2$, $M$ is diagonalisable and $u_n=\alpha r_1^n+\beta r_2^n$ for constants $\alpha,\beta$ fixed by $u_0,u_1$.
\end{example}

\section{Exercises}

\begin{exercise}[$\star$]\label{ex:fibmat}
Let $M=\begin{pmatrix}1&1\\1&0\end{pmatrix}$ and $(F_n)$ the Fibonacci numbers.
\begin{enumerate}[label=(\alph*)]
\item Show that $M^n=\begin{pmatrix}F_{n+1}&F_n\\F_n&F_{n-1}\end{pmatrix}$ for $n\ge1$, and deduce Cassini's identity $F_{n+1}F_{n-1}-F_n^2=(-1)^n$.
\item Diagonalise $M$ and deduce Binet's formula $F_n=\frac{\varphi^n-\psi^n}{\sqrt5}$, where $\varphi=\frac{1+\sqrt5}2$ and $\psi=\frac{1-\sqrt5}2$.
\end{enumerate}
\end{exercise}

\begin{exercise}[$\star$]\label{ex:nilp}
Let $A\in M_n(\C)$ be nilpotent. Show that $\operatorname{tr}A=0$ and $\det(I_n+A)=1$.
\end{exercise}

\begin{exercise}[$\star$]\label{ex:mindet}
Compute the determinant of the $n\times n$ matrix $\left(\min(i,j)\right)_{1\le i,j\le n}$.
\end{exercise}

\begin{exercise}[$\star\star$]\label{ex:annul}
Let $A\in M_n(\R)$ with $A^2-3A+2I_n=0$. Show that $A$ is invertible, compute $A^{-1}$, and show that $A$ is diagonalisable. What are the possible eigenvalues?
\end{exercise}

\begin{exercise}[$\star\star$]\label{ex:abba}
Let $A,B\in M_n(\C)$. Show that $AB$ and $BA$ have the same characteristic polynomial. \emph{Hint: start with $A$ invertible.}
\end{exercise}

\begin{exercise}[$\star\star$]\label{ex:rayleigh}
Let $S$ be real symmetric with eigenvalues $\lambda_1\le\dots\le\lambda_n$.
\begin{enumerate}[label=(\alph*)]
\item Show that $\max_{\|x\|=1}x^{\mathsf T}Sx=\lambda_n$ and $\min_{\|x\|=1}x^{\mathsf T}Sx=\lambda_1$.
\item Show that if $x^{\mathsf T}Sx\ge0$ for all $x$, there is a symmetric $R$ with $x^{\mathsf T}Rx\ge0$ for all $x$ and $R^2=S$.
\end{enumerate}
\end{exercise}

\begin{exercise}[$\star\star\star$]\label{ex:hilbert}
\emph{The Hilbert matrix.} Let $H_n=\left(\frac1{i+j-1}\right)_{1\le i,j\le n}$. Show that $H_n$ is the Gram matrix of $1,x,\dots,x^{n-1}$ for the inner product
$\langle f,g\rangle=\int_0^1fg$, and deduce that $\det H_n>0$.
\end{exercise}

\begin{exercise}[$\star\star\star$]\label{ex:codiag}
Let $A,B\in M_n(\C)$ be diagonalisable with $AB=BA$. Show that there is a single invertible $P$ such that $P^{-1}AP$ and $P^{-1}BP$ are both diagonal.
\end{exercise}

\begin{beauty}
Cassini's identity comes from $\det(M^n)=(\det M)^n$: an arithmetic fact read off a determinant. The spectral theorem is the engine of
Tome~3 (covariance matrices, Markov chains) and of PageRank in the notebook. Simultaneous diagonalisation (Exercise~\ref{ex:codiag}) is the
seed of the representation theory of Tome~2, where a whole commutative group of matrices is diagonalised at once by characters.
The Hilbert matrix is positive definite, yet its determinant is so small that computers cannot invert it: see the notebook.
\end{beauty}

\section{Solutions}

\begin{solution}{ex:fibmat}
(a) True for $n=1$ (with $F_0=0$), and $M^{n+1}=M^nM=\begin{pmatrix}F_{n+1}+F_n&F_{n+1}\\F_n+F_{n-1}&F_n\end{pmatrix}$ has the right form.
Taking determinants, $F_{n+1}F_{n-1}-F_n^2=(\det M)^n=(-1)^n$.

(b) $\chi_M=X^2-X-1$ has roots $\varphi\ne\psi$, so $M$ is diagonalisable: $M=PDP^{-1}$ and $M^n=PD^nP^{-1}$. Each entry of $M^n$ is therefore of the form
$\alpha\varphi^n+\beta\psi^n$; for $F_n$, the values $F_0=0$, $F_1=1$ give $\alpha+\beta=0$, $\alpha\varphi+\beta\psi=1$, so $\alpha=-\beta=\frac1{\varphi-\psi}=\frac1{\sqrt5}$.
\end{solution}

\begin{solution}{ex:nilp}
If $A^k=0$ and $Av=\lambda v$ with $v\ne0$, then $0=A^kv=\lambda^kv$, so $\lambda=0$. Over $\C$, $\chi_A$ splits, and all its roots are $0$: $\chi_A=X^n$. Hence
$\operatorname{tr}A=0$, and the eigenvalues of $I_n+A$ are all $1$ (as $\chi_{I+A}(X)=\chi_A(X-1)=(X-1)^n$), so $\det(I_n+A)=1$.
\end{solution}

\begin{solution}{ex:mindet}
For $i=n,n-1,\dots,2$ in that order, subtract row $i-1$ from row $i$. The new row $i$ has entries $\min(i,j)-\min(i-1,j)$, which is $1$ for $j\ge i$ and $0$ for $j<i$.
The matrix becomes upper triangular with ones on the diagonal: the determinant is $1$.
\end{solution}

\begin{solution}{ex:annul}
$A(3I_n-A)=2I_n$, so $A$ is invertible with $A^{-1}=\frac12(3I_n-A)$. The polynomial $X^2-3X+2=(X-1)(X-2)$ is split with simple roots and cancels $A$,
so $A$ is diagonalisable, with eigenvalues in $\{1,2\}$ (any eigenvalue $\lambda$ satisfies $\lambda^2-3\lambda+2=0$).
\end{solution}

\begin{solution}{ex:abba}
If $A$ is invertible, $BA=A^{-1}(AB)A$ is similar to $AB$, so they have the same characteristic polynomial. In general, $A+tI_n$ is invertible for all but finitely many $t\in\C$
(those $-t$ that are eigenvalues of $A$). For such $t$ and every $x$,
\[
\det\left(xI_n-(A+tI_n)B\right)=\det\left(xI_n-B(A+tI_n)\right).
\]
For fixed $x$, both sides are polynomials in $t$ that agree at infinitely many points, so they agree at $t=0$.
\end{solution}

\begin{solution}{ex:rayleigh}
(a) Write $S=PDP^{\mathsf T}$ and $y=P^{\mathsf T}x$; $P$ is orthogonal so $\|y\|=\|x\|=1$, and $x^{\mathsf T}Sx=y^{\mathsf T}Dy=\sum\lambda_iy_i^2$. This lies between $\lambda_1\sum y_i^2=\lambda_1$
and $\lambda_n$, and the bounds are attained at unit eigenvectors.

(b) By (a) all $\lambda_i\ge0$. Let $R=P\sqrt DP^{\mathsf T}$ where $\sqrt D=\operatorname{diag}(\sqrt{\lambda_i})$. Then $R$ is symmetric, $R^2=PDP^{\mathsf T}=S$, and its eigenvalues
$\sqrt{\lambda_i}\ge0$ give $x^{\mathsf T}Rx\ge0$ by (a).
\end{solution}

\begin{solution}{ex:hilbert}
$\langle x^{i-1},x^{j-1}\rangle=\int_0^1x^{i+j-2}\dd x=\frac1{i+j-1}$. For a Gram matrix $G=\left(\langle e_i,e_j\rangle\right)$ of linearly independent vectors and $c\in\R^n\setminus\{0\}$,
$c^{\mathsf T}Gc=\left\|\sum c_ie_i\right\|^2>0$, because $\sum c_ie_i\ne0$. So $G$ is symmetric with all eigenvalues positive (by Exercise~\ref{ex:rayleigh}), and $\det G$, their product, is positive.
The monomials are linearly independent, so $\det H_n>0$.
\end{solution}

\begin{solution}{ex:codiag}
Let $\lambda_1,\dots,\lambda_r$ be the distinct eigenvalues of $A$ and $E_k=\ker(A-\lambda_kI_n)$; since $A$ is diagonalisable, $\C^n=E_1\oplus\dots\oplus E_r$.
Each $E_k$ is stable under $B$: if $Av=\lambda_kv$ then $A(Bv)=BAv=\lambda_kBv$. A polynomial with simple roots that cancels $B$ also cancels its restriction $B_k$ to $E_k$,
so $B_k$ is diagonalisable: $E_k$ has a basis of eigenvectors of $B$, which are automatically eigenvectors of $A$ (for $\lambda_k$). Gathering these bases gives a basis of $\C^n$ made of common eigenvectors.
\end{solution}
